\SourceLeaf{403}%
\[
\left.\begin{aligned}
\frac{r-q}{p^*-q}&=\frac{(a-b)(d-c)(ab-cd)^2}{(a-d)(b-c)(ad-bc)^2}.pp'\\
\frac{r-q}{p^*-r}&=\frac{(a-b)(d-c)(ab-cd)^2}{(a-c)(b-d)(ac-bd)^2}.pp'\\
\frac{p^*-r}{p^*-q}&=\frac{(a-c)(b-d)(ac-bd)^2}{(a-d)(b-c)(ad-bc)^2}
\end{aligned}\right\}\ \text{sind reelle Zahlen}
\]
oder \(p^*,q,r\) liegen in einer geraden Linie normal gegen \(0p^*p\)
\par\medskip\Needspace{5\baselineskip}
\pdfbookmark[2]{V.}{carnot-zusatz-v}
\centerline{\textbf{V.}}\nobreak\smallskip
\{In einem gegebenen Kreise ein Vieleck zu beschreiben, dessen Seiten durch eben so viel gegebene Punkte gehen. \textsc{Schumacher}.\}

Es sei der Halbmesser des Kreises, \(r\), die Coordinaten der Winkelpunkte des Polygons
\[
\begin{array}{llll}
r\cos\varphi,&r\cos\varphi',&r\cos\varphi''&\text{etc.}\\
r\sin\varphi,&r\sin\varphi',&r\sin\varphi''&\text{etc.}
\end{array}
\]
endlich die Coordinaten der gegebenen Punkte, durch welche die verlängerten Seiten des Polygons gehn, (welche respective den ersten und zweiten Winkelpunkt, den zweiten und dritten u.\,s.\,w. verbinden)
\[
\begin{array}{llll}
a\cos A,&a'\cos A',&a''\cos A''&\text{etc.}\\
a\sin A,&a'\sin A',&a''\sin A''&\text{etc.}
\end{array}
\]
\par
Dann ist nach dem Grundsatze, dass, wenn drei Punkte, deren Coordinaten \(x,y;\ x',y';\ x'',y''\) sind, in einer geraden Linie liegen, die Bedingungsgleichung
\[
xy'+x'y''+x''y-x'y-x''y'-xy''=0
\]
Statt hat,
\[
\left.\begin{aligned}
rr\cos\varphi\sin\varphi'+ra\cos\varphi'\sin A+ar\cos A\sin\varphi\\
-rr\cos\varphi'\sin\varphi-ar\cos A\sin\varphi'-ra\cos\varphi\sin A
\end{aligned}\right\}=0
\]
oder
\[
r\sin(\varphi'-\varphi)-a\sin(\varphi'-A)+a\sin(\varphi-A)=0
\]
oder
\[
r\cos\tfrac12(\varphi'-\varphi)=a\cos\tfrac12(\varphi'+\varphi-2A)
\]
